\chapter{The Distributed Runtime}
\label{ch:distributed-runtime}

\newthought{Under the hood.} Nothing in this chapter is needed to \emph{use} \Jtwo\ ---
the defaults run themselves --- but when you operate scans at scale, tune performance,
or debug something odd, a precise picture of the process family and its shared state
pays for itself. This is that picture.

\section{The process family}
\label{sec:process-family}

A running scan is five kinds of process:

\begin{table}[ht]
  \footnotesize
  \begin{tabular}{@{}p{0.32\linewidth}p{0.62\linewidth}@{}}
    \toprule
    \textbf{Process} & \textbf{Owns} \\
    \midrule
    \code{Jarvis2:<scan>} & the control process: the Sampler, the factory that spawns
      and supervises everything else, checkpointing, the run summary \\
    \code{Jarvis-Redis:<scan>} & the broker (started for you when not already running) \\
    \code{Jarvis2-Worker-N} & one Sample at a time: mapping, calculators, operas,
      likelihood; plus a per-Worker file-operation helper process for
      \yamlkey{save} copies and deletions \\
    \code{Jarvis2-Archiver} & batched database writes, bucket packing \\
    \bottomrule
  \end{tabular}
  \caption{The process family, as named in \clih{ps}.}
  \label{tab:process-family}
\end{table}

Workers are started with the \code{spawn} multiprocessing context: each is a fresh
interpreter that rebuilds its own state from the configuration --- imports, compiled
expressions, calculator templates, environment captures. Nothing live is inherited, so
Workers on any platform behave identically and nothing unpicklable ever crosses a
process boundary.

\section{What lives in Redis}
\label{sec:redis-schema}

Redis is the \emph{only} shared truth, and it holds only light data --- IDs, counters,
small dictionaries. Products live on disk.

\begin{table}[ht]
  \footnotesize
  \begingroup
  \setlength{\tabcolsep}{0pt}
  \begin{tabular}{@{}p{0.24\linewidth}p{0.20\linewidth}p{0.50\linewidth}@{}}
    \toprule
    \textbf{Key} & \textbf{Type} & \textbf{Role} \\
    \midrule
    \code{hep:task\_queue}          & list & Sampler pushes tasks; Workers block-pop \\
    \code{calc:free:<name>}        & list & calculator concurrency pool (slot ids) \\
    \code{hep:archive\_queue}       & list & finished-Sample handoff to the Archiver \\
    \code{hep:worker:status:<id>}  & hash & heartbeat, current Sample, held packs, \textsc{pid}s \\
    \code{hep:sample:stats}        & hash & running / completed / failed \\
    \code{hep:calculator:status}   & hash & free/busy per pool \\
    \code{hep:*:op\_count}          & str  & monotonic change counters per subsystem \\
    \code{hep:feedback}            & list & observables feedback (\sampler{AdaptiveBridson}) \\
    \bottomrule
  \end{tabular}
  \endgroup
  \caption{The Redis schema (illustrative key shapes).}
  \label{tab:redis-schema}
\end{table}

A task on the queue is a small \textsc{json} document --- \textsc{uuid}, u-coordinates,
the execution plan, optional carried parameters --- everything else is rebuilt in the
Worker (Section~\ref{sec:uspace}).

\paragraph{Read/write separation.} During normal operation only Workers and the Sampler
\emph{write} to Redis. The factory's monitoring thread only \emph{reads}, gated by the
\code{op\_count} counters: a subsystem is re-fetched only when its counter moved. That
is why monitoring is free (Section~\ref{sec:monitor-cli}).

\section{The Worker loop}
\label{sec:worker-loop}

Each Worker's life, after one-time initialization (Redis client, mapper, one calculator
instance per type with templates preloaded, imported operators, compiled expressions,
the subprocess scheduler, the file-operation helper):

\begin{plaincode}
loop:
  task <- blpop(hep:task_queue)
  sample <- rebuild from task; allocate bucket slot; open sample logger
  for each layer in execution_plan:
      run this layer's calculators concurrently
        (each: acquire pool slot -> prepare runtime -> inputs ->
         commands -> outputs -> release slot)
  run operas in order; evaluate likelihood
  submit result -> hep:archive_queue        (success or failure alike)
  close sample; bump op_counts; heartbeat
\end{plaincode}

One Sample at a time per Worker --- the \emph{only} intra-Sample concurrency is between
same-layer calculators, throttled by the global pools. Workers compete on the shared
queue, so load balances itself: a Worker stuck on a slow point simply takes fewer
points.

\section{Layer 1 vs Layer 2: the two I/O planes}
\label{sec:two-layers}

The design draws a hard line between two kinds of file traffic:

\begin{description}[style=nextline, leftmargin=1.2em]
  \item[Layer 1 --- inside the Sample.] A calculator chain's own files (write input,
        run, read output) are synchronous and per-Sample --- they \emph{are} the
        critical path. \yamlkey{save} copies are offloaded to the per-Worker
        file-operation process so even they do not stall execution
        (Section~\ref{sec:save-policy}).
  \item[Layer 2 --- after the Sample.] Database rows, bucket packing, and any staging
        moves belong to the Archiver, batched (\yamlkey{archiver.batch\_size}, flush
        interval) and entirely off the Workers' path. By default the handoff is
        \emph{direct}: products are already in their final bucket location, and only
        records move.
\end{description}

Bucket packing is gated three ways --- sealed, no active Samples, all records written
--- so a \file{.tar.gz} can never contain a Sample whose database row does not exist
(Section~\ref{sec:sample-directory}).

\section{Failure semantics, precisely}
\label{sec:runtime-failures}

\begin{itemize}
  \item \textbf{Sample fails} (calculator non-zero exit, timeout, missing observable,
        opera \code{TypeError}): the Worker records a failed result --- archived like
        any other --- keeps the Sample's log and directory, releases slots, and
        continues.
  \item \textbf{Worker dies}: heartbeats lapse; the factory respawns a fresh Worker;
        the in-flight Sample is re-queued. Held pool slots are reclaimed.
  \item \textbf{Archiver dies}: results wait in the archive queue; the factory
        restarts it; batched writes are idempotent per \textsc{uuid}, so replays cannot
        duplicate rows.
  \item \textbf{Redis dies}: the scan cannot continue --- \Jtwo\ fails fast with a
        connection error rather than degrading. The checkpoint covers you: restart the
        service and \cli{Jarvis2 run --resume}.
  \item \textbf{Control process dies}: the 30-second checkpoint plus resume
        (Chapter~\ref{ch:checkpoint}).
\end{itemize}

\begin{jnote}
The recurring theme: \emph{Redis holds nothing precious}. Queues drain and rebuild;
pools reseed from the card; the Sampler's checkpoint plus the Archiver's acknowledged
records are the durable truth. That is what makes both crash recovery and resume exact.
\end{jnote}
